\documentclass{beamer}
\usepackage{amsmath, amssymb}
\usetheme{Madrid}
\title{GED Mathematics: Geometry Unit}
\author{MacLachlan College}
\date{\today}

\begin{document}

\frame{\titlepage}

% ------------------------------
% Lesson 1: Points, Lines, Angles
% ------------------------------
\begin{frame}{Lesson 1: Points, Lines, and Angles}
\begin{itemize}
  \item Points represent exact locations in space.
  \item Lines extend infinitely in two directions.
  \item Line segments have two endpoints.
  \item Rays start at one point and extend infinitely in one direction.
  \item Angles are formed by two rays sharing a common endpoint.
\end{itemize}
\end{frame}

\begin{frame}{Lesson 1: Worked Examples}
\textbf{Example 1:} Find the complement of a 35$^\circ$ angle.\\
Complement = 90$^\circ$ - 35$^\circ$ = 55$^\circ$.\\[1em]
\textbf{Example 2:} Find the supplement of a 120$^\circ$ angle.\\
Supplement = 180$^\circ$ - 120$^\circ$ = 60$^\circ$.\\[1em]
\textbf{Example 3:} If two angles form a linear pair and one is 70$^\circ$, the other is 110$^\circ$.\\[1em]
\textbf{Example 4:} If two complementary angles are in the ratio 1:2, find them.\\
$x + 2x = 90 \Rightarrow 3x = 90 \Rightarrow x = 30$. So angles are 30$^\circ$ and 60$^\circ$.\\[1em]
\textbf{Example 5:} Identify adjacent vs. vertical angles: Vertical angles are opposite each other when two lines intersect.
\end{frame}

% ------------------------------
% Lesson 2: Triangles
% ------------------------------
\begin{frame}{Lesson 2: Triangles}
\begin{itemize}
  \item The sum of angles in a triangle is 180$^\circ$.
  \item Types: scalene, isosceles, equilateral.
  \item Pythagorean theorem: $a^2 + b^2 = c^2$.
  \item Area formula: $A = \tfrac{1}{2}bh$.
  \item Perimeter is the sum of side lengths.
\end{itemize}
\end{frame}

\begin{frame}{Lesson 2: Worked Examples}
\textbf{Example 1:} Find the missing angle in a triangle with angles 70$^\circ$ and 50$^\circ$.\\
$180 - (70 + 50) = 60^\circ$.\\[1em]
\textbf{Example 2:} A right triangle has legs 3 and 4. Find the hypotenuse.\\
$c = \sqrt{3^2 + 4^2} = 5$.\\[1em]
\textbf{Example 3:} Find the area of a triangle with base 10 and height 6.\\
$A = \tfrac{1}{2}(10)(6) = 30$.\\[1em]
\textbf{Example 4:} For an isosceles triangle with equal sides 5 and base 6, find the height.\\
Height = $\sqrt{5^2 - 3^2} = 4$.\\[1em]
\textbf{Example 5:} Find the perimeter of a triangle with sides 5, 7, and 9.\\
$5 + 7 + 9 = 21$.
\end{frame}

% ------------------------------
% Lesson 3: Quadrilaterals
% ------------------------------
\begin{frame}{Lesson 3: Quadrilaterals}
\begin{itemize}
  \item Quadrilaterals have 4 sides and 360$^\circ$ in total.
  \item Types: square, rectangle, parallelogram, trapezoid, rhombus.
  \item Area formulas:
    \begin{itemize}
      \item Rectangle: $A = lw$
      \item Parallelogram: $A = bh$
      \item Trapezoid: $A = \tfrac{1}{2}(b_1 + b_2)h$
    \end{itemize}
\end{itemize}
\end{frame}

\begin{frame}{Lesson 3: Worked Examples}
\textbf{Example 1:} Find the area of a rectangle with $l=8, w=3$.\\
$A = 8 \times 3 = 24$.\\[1em]
\textbf{Example 2:} Area of a parallelogram with base 12 and height 5: $A = 60$.\\[1em]
\textbf{Example 3:} Area of a trapezoid with bases 10, 6 and height 4: $A = \tfrac{1}{2}(10+6)(4) = 32$.\\[1em]
\textbf{Example 4:} A square has diagonal 10. Find side length: $s = \tfrac{10}{\sqrt{2}} = 5\sqrt{2}$.\\[1em]
\textbf{Example 5:} Perimeter of rhombus with side 7: $4 \times 7 = 28$.
\end{frame}

% ------------------------------
% Lesson 4: Circles
% ------------------------------
\begin{frame}{Lesson 4: Circles}
\begin{itemize}
  \item Radius: distance from center to any point on circle.
  \item Diameter: twice the radius.
  \item Circumference: $C = 2\pi r$.
  \item Area: $A = \pi r^2$.
\end{itemize}
\end{frame}

\begin{frame}{Lesson 4: Worked Examples}
\textbf{Example 1:} Find circumference with $r=7$. $C = 2\pi(7) = 14\pi$.\\[1em]
\textbf{Example 2:} Find area with $r=5$. $A = 25\pi$.\\[1em]
\textbf{Example 3:} If circumference = $31.4$, find $r$. $r = \tfrac{31.4}{2\pi} = 5$.\\[1em]
\textbf{Example 4:} A circle has area $78.5$. Find $r$. $r = \sqrt{\tfrac{78.5}{\pi}} \approx 5$.\\[1em]
\textbf{Example 5:} Diameter = 10. Find circumference: $C = \pi d = 10\pi$.
\end{frame}

% ------------------------------
% Lesson 5: 3D Figures
% ------------------------------
\begin{frame}{Lesson 5: 3D Figures}
\begin{itemize}
  \item Common solids: prism, cylinder, pyramid, cone, sphere.
  \item Volume formulas:
  \begin{itemize}
    \item Prism: $V = Bh$
    \item Cylinder: $V = \pi r^2h$
    \item Pyramid: $V = \tfrac{1}{3}Bh$
    \item Cone: $V = \tfrac{1}{3}\pi r^2h$
    \item Sphere: $V = \tfrac{4}{3}\pi r^3$
  \end{itemize}
\end{itemize}
\end{frame}

\begin{frame}{Lesson 5: Worked Examples}
\textbf{Example 1:} Find volume of a rectangular prism with $l=4, w=3, h=5$.\\
$V = lwh = 60$.\\[1em]
\textbf{Example 2:} Cylinder with $r=3, h=10$: $V = 90\pi$.\\[1em]
\textbf{Example 3:} Cone with $r=4, h=9$: $V = \tfrac{1}{3}\pi(16)(9) = 48\pi$.\\[1em]
\textbf{Example 4:} Sphere with $r=6$: $V = \tfrac{4}{3}\pi(216) = 288\pi$.\\[1em]
\textbf{Example 5:} Pyramid with base area 30 and height 12: $V = 120$.
\end{frame}

\end{document}
